\label{ch:ew}

\section{Fonctions constitutive et déterminantale de l'opérateur angulaire}

Rappelons l'opérateur à deux feuillets du \cref{ch:vacuum},
\begin{equation}
\label{eq:ew-angular-operator}
T=q_+P_+ +q_-P_-,
\qquad
q_-=\ee^{-(A_{\rm rad}-C^*_{\rm rad})},
\qquad
q_+=\ee^{-(A_{\rm rad}+C^*_{\rm rad})}.
\end{equation}
Le vide électromagnétique applique la fonction
\begin{equation}
\label{eq:ew-vacuum-reader}
F(T)^2,
\qquad
F(z)=\frac{1-z^{90}}{1-z^{120}},
\end{equation}
tandis que le secteur angulaire électrofaible applique le déterminant
\begin{equation}
\label{eq:ew-determinant-reader}
\boxed{
\det T=q_+q_-=\ee^{-2A_{\rm rad}}=:D_A.}
\end{equation}
Ce sont deux réalisations fonctionnelles de l'antécédent commun \((A,C^*,T)\). Il n'existe pas de
flèche causale \(F(T)^2\to\det T\) ni son inverse. L'unité structurelle
réside dans \(T\), et non dans l'identification ultérieure de la perméabilité, de la permittivité et de
\(\theta_W\).

Dans l'involution des feuillets \(C^*\mapsto-C^*\), les valeurs propres
\(q_+,q_-\) s'échangent. C'est pourquoi \(D_A\) et la vitesse constitutive
sont paires, tandis que \((\widehat\mu,\widehat\varepsilon)\) s'échange et
\(\widehat Z\) s'inverse. La relation électrofaible qui suit utilise
l'invariant pair.

\section{Relation exacte entre les secteurs \(W\) et \(Z\)}

La coordonnée angulaire de durée--perduration est
\begin{equation}
\label{eq:ew-angular-character}
\chi_{\rm dur}(n_A,s_C)
=\exp\!\left(n_AA_{\rm rad}
+\frac{s_C}{6}C^*_{\rm rad}\right).
\end{equation}
Cette coordonnée n'est pas introduite comme une étiquette isolée. Sa provenance est
\begin{equation}
\label{eq:ew-signature-provenance}
\begin{aligned}
\APP&\longrightarrow\TRIT\longrightarrow\TPK
\longrightarrow\text{extension survivante}
\longrightarrow\text{route}\\
&\longrightarrow\text{registre}
\longrightarrow(n_A,s_C,\ldots)
\longrightarrow\chi_{\rm dur}.
\end{aligned}
\end{equation}
Les signatures de \(W\), \(Z\) et \(H\) sont des grandeurs typées obtenues au moyen de
cette chaîne ; le présent chapitre les adopte sans les sélectionner à partir de
leurs masses terminales. Dans la réalisation angulaire établie, les secteurs
chargé et neutre satisfont
\begin{equation}
\label{eq:ew-WZ-signatures}
(n_A,s_C)_W=(94,-1),
\qquad
(n_A,s_C)_Z=(95,-1).
\end{equation}
Cette affirmation concerne la composante angulaire commune. Les
coordonnées supplémentaires de raffinement de la loi générale des masses ne sont pas
éliminées à moins que le transducteur physique ne démontre leur annulation.

\begin{theorem}[Défaut de Weinberg dans le schéma des masses physiques de la carte angulaire]
\label{thm:ew-weinberg}
Si \(W\) et \(Z\) sont évalués dans une même section d'échelle et partagent
le raffinement qui accompagne \cref{eq:ew-WZ-signatures}, alors
\begin{equation}
\label{eq:ew-WZ-ratio}
\boxed{
\frac{m_W^{\angle}}{m_Z^{\angle}}
=\ee^{-A_{\rm rad}}=\sqrt{D_A},
\qquad
\left(\frac{m_W^{\angle}}{m_Z^{\angle}}\right)^2=D_A.}
\end{equation}
Dans le schéma des masses physiques,
\begin{equation}
\label{eq:ew-sin2thetaW}
\boxed{
\sin^2\theta_W^{\rm OS,HMT}=1-D_A
=0.2248708807528435698111159035\ldots .}
\end{equation}
Adem\'as,
\begin{equation}
\label{eq:ew-WZ-terminal-ratio}
\frac{m_W^{\angle}}{m_Z^{\angle}}
=0.8804141748331613670128885459\ldots .
\end{equation}
\end{theorem}

\begin{proof}
Par \cref{eq:ew-angular-character,eq:ew-WZ-signatures},
\[
\frac{\chi_{\rm dur}(94,-1)}{\chi_{\rm dur}(95,-1)}
=\exp(-A_{\rm rad});
\]
la coordonnée \(C^*\) s'annule parce que les deux secteurs partagent
\(s_C=-1\). En élevant au carré et en utilisant
\cref{eq:ew-determinant-reader}, on obtient \(D_A\). La définition
des masses physiques \(s_W^2=1-m_W^2/m_Z^2\) produit
\cref{eq:ew-sin2thetaW}.
\end{proof}

Le contenu principal est l'identité entre deux constructions
internes : le déterminant de l'opérateur angulaire et la différence d'un pas
dans la réalisation \(W/Z\). La promotion en sélecteur physique de jauge exige
un entrelaceur qui transporte les feuillets de \(T\) vers la base neutre
\((W^3,B)\) et qui démontre quels raffinements sont conservés ou annulés.

\section{Couplages de jauge en normalisation rationalisée}

Nous définissons le couplage adimensionnel de jauge
\begin{equation}
\label{eq:ew-egauge}
e_g:=\sqrt{4\pi\alpha_{\HMT}}
=0.3028221208717526427662442689\ldots .
\end{equation}
Il ne faut pas confondre \(e_g\) avec une section dimensionnelle de charge
\(e_{\rm int}\) : la première est un couplage en unités naturelles ; la
seconde vit dans la droite de charge de la mécanique dimensionnelle.

Soit
\begin{equation}
\label{eq:ew-sw-cw}
s_W^2=1-D_A,
\qquad
c_W^2=D_A.
\end{equation}
La convention électrofaible rationalisée au niveau de l'arbre donne
\begin{align}
g&=\frac{e_g}{s_W}
=2\sqrt{\frac{\pi\alpha_{\HMT}}{1-D_A}}
=0.6385883427334574283647003809\ldots,
\label{eq:ew-g}\\
g'&=\frac{e_g}{c_W}
=2\sqrt{\frac{\pi\alpha_{\HMT}}{D_A}}
=0.3439541633108502429362319257\ldots,
\label{eq:ew-gprime}\\
\frac{g'}{g}
&=\sqrt{D_A^{-1}-1}
=0.5386164141966093594996610933\ldots .
\label{eq:ew-gprime-over-g}
\end{align}

\begin{proposition}[Relation custodiale à l'arbre]
\label{prop:ew-rho}
Dans la même interface,
\begin{equation}
\label{eq:ew-rho}
\rho_{\rm EW}^{(0)}
:=\frac{(m_W^{\angle})^2}
{(m_Z^{\angle})^2c_W^2}=1.
\end{equation}
\end{proposition}

\begin{proof}
Par \cref{eq:ew-WZ-ratio,eq:ew-sw-cw}, le numérateur relatif comme
\(c_W^2\) sont \(D_A\).
\end{proof}

Les équations \cref{eq:ew-g,eq:ew-gprime} sont des compositions exactes une
fois déclarés le schéma des masses physiques, la rationalisation et l'échelle. Elles
n'affirment pas qu'un changement de schéma de renormalisation laisse leurs
valeurs numériques inchangées.

\section{Constante de Fermi dans l'approximation à l'arbre}

En unités naturelles, l'intégration du propagateur chargé à basse
énergie produit
\begin{equation}
\label{eq:ew-GF-def}
\frac{G_F^{(0)}}{\sqrt2}=\frac{g^2}{8m_W^2}.
\end{equation}
En substituant \cref{eq:ew-g}, on obtient la composition
\begin{equation}
\label{eq:ew-GF}
\boxed{
G_F^{(0)}
=\frac{\pi\alpha_{\HMT}}
{\sqrt2(1-D_A)m_W^2}.}
\end{equation}
En adoptant la masse obtenue précédemment
\begin{equation}
\label{eq:ew-W-output}
m_W^{\HMT}=80.381592550221\ldots\ \mathrm{GeV},
\end{equation}
l'évaluation terminale est
\begin{equation}
\label{eq:ew-GF-value}
\boxed{
G_F^{(0)}
=1.1157162817659203186621784\times10^{-5}\ \mathrm{GeV}^{-2}.}
\end{equation}

La différence par rapport à une constante de Fermi renormalisée n'est pas utilisée
pour choisir un coefficient. Dans la convention
\begin{equation}
\label{eq:ew-Delta-r}
G_F
=\frac{\pi\alpha_{\HMT}}
{\sqrt2(1-D_A)m_W^2(1-\Delta r)},
\end{equation}
l'obligation restante est de dériver \(\Delta r\) à partir du secteur interne des
boucles, des états et de l'échelle. Consulter la valeur observée de \(G_F\) pour définir
\(\Delta r\) serait une reconstruction calibrée, et non une prédiction.

\section{Dépendance orientée du couplage de Higgs}

Dans la carte angulaire, les routes pertinentes sont
\begin{equation}
\label{eq:ew-HW-signatures}
(n_A,s_C)_H=(97,9),
\qquad
(n_A,s_C)_W=(94,-1).
\end{equation}
Par conséquent,
\begin{equation}
\label{eq:ew-HW-ratio}
\boxed{
\frac{m_H^{\angle}}{m_W^{\angle}}
=\frac{\chi_{\rm dur}(97,9)}{\chi_{\rm dur}(94,-1)}
=\exp\!\left(3A_{\rm rad}+\frac53C^*_{\rm rad}\right)
=1.55872087212325548564\ldots .}
\end{equation}
Ici, \(C^*\) ne s'annule pas : le secteur scalaire conserve l'orientation
que le déterminant \(D_A=\ee^{-2A_{\rm rad}}\) ne distingue pas.

Prenons la convention du potentiel
\begin{equation}
\label{eq:ew-Higgs-potential}
V(H)=-\mu_H^2H^\dagger H
+\lambda_H(H^\dagger H)^2,
\end{equation}
pour laquelle
\begin{equation}
\label{eq:ew-Higgs-tree-relations}
m_H^2=2\lambda_Hv^2,
\qquad
m_W^2=\frac{g^2v^2}{4}.
\end{equation}

\begin{corollary}[Autocouplage orienté de Higgs]
\label{cor:ew-Higgs-lambda}
Dans l'interface à l'arbre déclarée,
\begin{equation}
\label{eq:ew-Higgs-lambda}
\boxed{
\lambda_H
=\frac{\pi\alpha_{\HMT}}{2(1-D_A)}
\exp\!\left(6A_{\rm rad}+\frac{10}{3}C^*_{\rm rad}\right)
=0.1238479115482466804662\ldots .}
\end{equation}
Adem\'as,
\begin{equation}
\label{eq:ew-GF-Higgs}
G_F^{(0)}m_H^2=\sqrt2\,\lambda_H.
\end{equation}
\end{corollary}

\begin{proof}
Diviser les deux équations de
\cref{eq:ew-Higgs-tree-relations} donne
\(m_H^2/m_W^2=8\lambda_H/g^2\). Avec
\(g^2=4\pi\alpha_{\HMT}/(1-D_A)\) et le carré de
\cref{eq:ew-HW-ratio}, on obtient \cref{eq:ew-Higgs-lambda}. D'autre part,
\(G_F^{(0)}=1/(\sqrt2v^2)\) et
\(m_H^2=2\lambda_Hv^2\) impliquent
\cref{eq:ew-GF-Higgs}.
\end{proof}

La valeur de \(\lambda_H\) est antérieure à l'incorporation des corrections
radiatives et dépend de
la convention et de l'échelle déclarées. Son intérêt structurel est qu'elle
sépare deux mémoires : le secteur commun de jauge emploie l'invariant pair \(D_A\),
tandis que la route de Higgs retient explicitement \(C^*\).

\section{Clôture déterminantale du secteur électrofaible : mélange, couplages et spectre constitutif}

La chaîne complète de ce chapitre est
\begin{equation}
\label{eq:ew-causal-diagram}
\boxed{
\begin{aligned}
(A,C^*)&\longrightarrow T
\longrightarrow
\begin{cases}
F(T)^2&\longrightarrow
(\widehat\mu,\widehat\varepsilon,\widehat Z,\widehat c),\\
\det T=D_A&\longrightarrow
(s_W^2,c_W^2),
\end{cases}\\
(D_A,\alpha_{\HMT})&\longrightarrow(e_g,g,g'),\\
(D_A,\alpha_{\HMT},m_W)&\longrightarrow G_F^{(0)},\\
(A,C^*,D_A,\alpha_{\HMT})&\longrightarrow
\left(\frac{m_H}{m_W},\lambda_H\right).
\end{aligned}}
\end{equation}
Le chiffre apparaît toujours à la fin. Ni \(\sin^2\theta_W\), ni
\(G_F\), ni \(\lambda_H\) ne sont introduits comme valeurs cibles pour sélectionner
\(A,C^*\) ou les signatures de route.

\begin{proposition}[Résultat principal]
Le même opérateur angulaire détermine le vide par une fonction rationnelle
de ses feuillets et le défaut électrofaible par son déterminant. La
différence d'un pas entre \(W\) et \(Z\) transforme ce déterminant en
\(\sin^2\theta_W^{\rm OS}=1-D_A\) ; la route de Higgs, en revanche, conserve
le contre-angle et conserve l'orientation éliminée par le déterminant.
\end{proposition}

\begin{remark}[Contrôle négatif]
La relation de Weinberg est réfutée dans son domaine si les signatures
angulaires \(W/Z\) ne diffèrent pas uniquement par \(n_A:94\to95\), si une coordonnée
de raffinement ne s'annule pas dans l'interface déclarée ou si
\((m_W^{\angle}/m_Z^{\angle})^2\neq D_A\). La promotion de jauge est
bloquée s'il n'existe pas d'entrelaceur avec la base \((W^3,B)\). Les chiffres
de \(g,g',G_F,\lambda_H\) changent si l'on modifie le schéma ou l'échelle sans
l'enregistrer. Utiliser les valeurs observées de \(G_F\) ou \(m_H\) pour sélectionner
\(\Delta r\), les signatures ou \(C^*\) introduit une rétroaction par la
valeur cible et invalide la dérivation. Confondre \(e_g\) avec la charge
dimensionnelle constitue une erreur de type.
\end{remark}

\subsection*{Portée logique et domaine de validité}
Déterminant angulaire, quotient \(W/Z\) et défaut complémentaire de Weinberg :
\emph{Théorème interne exact dans la carte angulaire conditionnelle}. Couplages \(g,g'\),
\(\rho_{\rm EW}^{(0)}\), Fermi et Higgs : \emph{Composition exacte dans
l'interface électrofaible à l'arbre}. L'entrelaceur de base de jauge et la
correction interne \(\Delta r\) : \emph{Frontière ouverte avec critère de réfutation}.
Fermi et Higgs ne sont pas des constantes indépendantes sélectionnées sans valeurs
cibles : ce sont des compositions exactes qui emploient les signatures, la grandeur
\(m_W\) et la
convention à l'arbre déclarée. Les valeurs conventionnelles ultérieures
sont une \emph{Comparaison métrologique externe}.

